\chapter{Дифференциальные уравнения, динамика и дифференцируемые модели}
\label{chap:dynamics-differentiable-models}

Во многих задачах машинного обучения объект задаётся не одним статическим вектором, а процессом: конструкция деформируется под нагрузкой, температура распределяется во времени, скрытое состояние модели последовательно преобразуется, а параметры проектирования влияют на всю траекторию. Дифференциальное уравнение записывает такой процесс через закон мгновенного изменения состояния.

Для генеративного проектирования важна не коллекция специальных приёмов интегрирования, а вычислительная цепочка
\[
\begin{aligned}
\text{параметры конструкции}
&\longrightarrow
\text{динамическая модель}
\longrightarrow
\text{траектория состояния},\\
\text{траектория состояния}
&\longrightarrow
\text{целевая характеристика}
\longrightarrow
\text{градиент}.
\end{aligned}
\]
Эта глава вводит ровно тот аппарат, который нужен для понимания такой цепочки: корректность задачи Коши, линейная и нелинейная динамика, устойчивость, численное интегрирование и дифференцирование решения по параметрам.

\section{Состояние системы и корректность задачи Коши}
\label{sec:state-and-well-posed-ivp}

\subsection{От статической зависимости к закону эволюции}

Статическая модель сопоставляет входу $u$ один выход $y=F(u)$. Динамическая модель вместо одного выхода строит функцию времени
\[
x:I\to\R^n,
\qquad
 t\mapsto x(t),
\]
где $x(t)$ содержит всю информацию, необходимую для продолжения процесса после момента $t$. Такой вектор называется \emph{состоянием}, а функция $x(\cdot)$ --- \emph{траекторией состояния}.

\begin{definition}
Обыкновенным дифференциальным уравнением первого порядка называется соотношение
\begin{equation}
\label{eq:general-first-order-ode}
\dot x(t)=f\bigl(t,x(t)\bigr),
\qquad
x(t)\in\R^n,
\end{equation}
где $f:U\subseteq\R\times\R^n\to\R^n$ задаёт мгновенную скорость изменения состояния.
\end{definition}

Правая часть $f$ называется векторным полем или законом динамики. Если координата $x_i$ измеряется в единицах $[x_i]$, то координата $f_i$ имеет размерность $[x_i]/[t]$. Проверка размерностей часто выявляет ошибку модели раньше, чем вычисления.

Если $f$ не зависит от времени явно,
\begin{equation}
\label{eq:autonomous-ode-v01}
\dot x=f(x),
\end{equation}
система называется автономной. В ней одинаковое состояние всегда имеет одинаковую мгновенную скорость. Неавтономную систему можно формально сделать автономной, добавив время к состоянию:
\[
z=(t,x),
\qquad
\dot z=
\begin{pmatrix}
1\\
f(t,x)
\end{pmatrix}.
\]
Поэтому система первого порядка является основной формой ОДУ.

Уравнение порядка $k$
\[
x^{(k)}
=
F\bigl(t,x,\dot x,\ldots,x^{(k-1)}\bigr)
\]
сводится к ней заменой
\[
y_1=x,
\quad
y_2=\dot x,
\quad\ldots\quad,
y_k=x^{(k-1)}.
\]
Тогда
\begin{equation}
\label{eq:higher-order-to-first-order-v01}
\dot y_1=y_2,
\quad
\dot y_2=y_3,
\quad\ldots\quad,
\dot y_k=F(t,y_1,\\ldots,y_k).
\end{equation}

\begin{application}[динамическая модель конструкции]
Пусть $q(t)\in\R^d$ описывает перемещения узлов или обобщённые координаты конструкции, а $\eta$ --- геометрические и материальные параметры проекта. Модель второго порядка
\begin{equation}
\label{eq:mechanical-design-second-order-v01}
M(\eta)\ddot q
+C(\eta)\dot q
+K(\eta)q
=p(t)
\end{equation}
становится системой первого порядка для состояния $x=(q,v)$:
\begin{equation}
\label{eq:mechanical-design-first-order-v01}
\dot x
=
\begin{pmatrix}
\dot q\\
\dot v
\end{pmatrix}
=
\begin{pmatrix}
v\\
M(\eta)^{-1}
\bigl(p(t)-C(\eta)v-K(\eta)q\bigr)
\end{pmatrix}.
\end{equation}
Так одна точка $x(t)$ содержит и положение, и скорость; именно поэтому она определяет дальнейшее движение. Параметры $\eta$ влияют не на один ответ, а на всё векторное поле и всю траекторию.
\end{application}

\subsection{Задача Коши и интегральная форма}

Одного закона динамики недостаточно: через разные начальные состояния проходят разные траектории.

\begin{definition}
Задачей Коши называется система
\begin{equation}
\label{eq:ivp-v01}
\dot x(t)=f\bigl(t,x(t)\bigr),
\qquad
x(t_0)=x_0.
\end{equation}
Решением на интервале $I\ni t_0$ называется дифференцируемая функция $x:I\to\R^n$, удовлетворяющая обоим равенствам.
\end{definition}

Интегрирование от $t_0$ до $t$ даёт эквивалентную форму
\begin{equation}
\label{eq:ivp-integral-form-v01}
x(t)
=
x_0+
\int_{t_0}^{t}
f\bigl(s,x(s)\bigr)\,\dd s.
\end{equation}
Обратно, если непрерывная функция удовлетворяет~\eqref{eq:ivp-integral-form-v01}, то по основной теореме анализа она дифференцируема и решает~\eqref{eq:ivp-v01}. Интегральная запись важна потому, что превращает поиск траектории в задачу о неподвижной точке оператора
\[
(\mathcal Kx)(t)
=
x_0+
\int_{t_0}^{t}f\bigl(s,x(s)\bigr)\,\dd s.
\]

Начав с постоянной функции $x^{(0)}(t)=x_0$, можно строить последовательность
\begin{equation}
\label{eq:picard-iteration-v01}
x^{(m+1)}(t)
=
x_0+
\int_{t_0}^{t}
f\bigl(s,x^{(m)}(s)\bigr)\,\dd s.
\end{equation}
Это итерации Пикара. При подходящих условиях они сходятся к единственному решению. Здесь важна не практическая эффективность схемы, а идея: существование траектории выводится из сходимости последовательных уточнений.

\subsection{Почему непрерывности недостаточно}

Непрерывность $f$ гарантирует локальное существование хотя бы одного решения, но сама по себе не запрещает нескольким траекториям выходить из одной точки.

\begin{definition}
Функция $f(t,x)$ называется локально липшицевой по состоянию $x$, если для каждой точки $(t_0,x_0)$ существует окрестность и число $L\geq0$, для которых
\begin{equation}
\label{eq:local-lipschitz-v01}
\norm{f(t,x)-f(t,y)}
\leq
L\norm{x-y}
\end{equation}
при всех допустимых $t,x,y$ из этой окрестности.
\end{definition}

Условие ограничивает скорость изменения векторного поля по состоянию. Если $f$ непрерывно дифференцируема по $x$, то на каждом компактном рабочем множестве можно взять
\[
L
\geq
\sup_{(t,x)}\norm{D_xf(t,x)}.
\]
Следовательно, гладкие модели локально липшицевы. Кусочно-линейные нейронные сети с конечными весами также липшицевы; грубую верхнюю оценку константы даёт произведение операторных норм весовых матриц и констант Липшица активаций.

\begin{theorem}[локальное существование и единственность]
\label{thm:picard-lindelof-v01}
Пусть $f$ непрерывна в окрестности $(t_0,x_0)$ и локально липшицева по $x$. Тогда существует интервал $(t_0-\tau,t_0+\tau)$, на котором задача Коши~\eqref{eq:ivp-v01} имеет единственное решение.
\end{theorem}

\begin{proofidea}
На достаточно коротком интервале оператор $\mathcal K$ сохраняет замкнутый шар непрерывных траекторий и становится сжатием в равномерной норме. Принцип сжимающих отображений даёт единственную неподвижную точку, а итерации~\eqref{eq:picard-iteration-v01} сходятся к ней.
\end{proofidea}

Требование Липшица существенно именно для единственности. Рассмотрим
\begin{equation}
\label{eq:nonunique-sqrt-v01}
\dot x=\sqrt{\abs{x}},
\qquad
x(0)=0.
\end{equation}
Правая часть непрерывна, но не липшицева в нуле. Постоянная функция $x(t)=0$ является решением. Для любого $a\geq0$ решением также служит
\[
x_a(t)
=
\begin{cases}
0, & t\leq a,\\[2mm]
\dfrac{(t-a)^2}{4}, & t>a.
\end{cases}
\]
Система может произвольно долго оставаться в нуле, а затем начать движение. Поэтому одно начальное состояние не задаёт однозначный выход.

\begin{mlconnection}[детерминированность непрерывного слоя]
Пусть скрытое состояние модели удовлетворяет
\[
\dot z=f_\theta(t,z),
\qquad
z(0)=E_\theta(u),
\]
а предсказание вычисляется как $\widehat y=H_\theta(z(T))$. При единственности решение задаёт функцию
\[
(u,\theta)\longmapsto z(T)\longmapsto\widehat y.
\]
Без единственности одна и та же пара $(u,\theta)$ допускает несколько конечных состояний, поэтому модель не определяет однозначного прямого прохода. Условие локальной липшицевости является математической формой требования, что непрерывный слой действительно задаёт отображение.
\end{mlconnection}

\subsection{Локальное решение не обязательно глобально}

Теорема~\ref{thm:picard-lindelof-v01} гарантирует решение лишь вблизи $t_0$. Его можно продолжать, пока траектория остаётся внутри области определения и не уходит на бесконечность. Максимальным называется решение, продолжённое на наибольший возможный интервал.

Даже гладкое поле может вызвать разрушение за конечное время. Для задачи
\begin{equation}
\label{eq:finite-time-blow-up-v01}
\dot x=x^2,
\qquad
x(0)=x_0>0,
\end{equation}
получаем
\begin{equation}
\label{eq:finite-time-blow-up-solution-v01}
x(t)=\frac{x_0}{1-x_0t}.
\end{equation}
Решение единственно, но определено только при $t<1/x_0$ и неограниченно возрастает при $t\uparrow1/x_0$.

\begin{caution}
Локальная корректность, глобальное существование, численная устойчивость и физическая адекватность --- разные свойства. Теорема существования не гарантирует, что решение существует до требуемого момента $T$, что его удастся точно вычислить или что выбранное уравнение правильно описывает объект.
\end{caution}

Для вычислительной модели обычно задают рабочую область $K$ состояний и интервал $[t_0,T]$. Если решение остаётся в компактном $K$, а $f$ непрерывна и локально липшицева в окрестности $[t_0,T]\times K$, то траекторию можно последовательно продолжить до $T$. На практике это связывает корректность с априорными ограничениями: допустимыми температурами, деформациями, концентрациями или нормами скрытого состояния.

\subsection{Зависимость от начальных данных и ошибки модели}

Единственность отвечает на вопрос, совпадают ли две траектории при одинаковых данных. Для применения важнее количественно оценить, насколько они расходятся при небольшом изменении данных.

Пусть $x$ и $y$ решают одно уравнение $\dot x=f(t,x)$, а $f$ имеет константу Липшица $L$ на общей рабочей области. Из интегральной формы следует
\[
\norm{x(t)-y(t)}
\leq
\norm{x_0-y_0}
+
L\int_{t_0}^{t}\norm{x(s)-y(s)}\,\dd s.
\]
Неравенство Гронуолла даёт оценку
\begin{equation}
\label{eq:initial-data-stability-v01}
\norm{x(t)-y(t)}
\leq
\exp\bigl(L(t-t_0)\bigr)
\norm{x_0-y_0},
\qquad t\geq t_0.
\end{equation}
Следовательно, на любом конечном интервале решение непрерывно зависит от начального состояния. Однако множитель $e^{L(t-t_0)}$ показывает, что корректность не исключает сильного усиления ошибки на длинном горизонте.

Пусть теперь $x$ решает $\dot x=f(t,x)$, а $\widetilde x$ --- приближённую модель
\[
\dot{\widetilde x}=\widetilde f(t,\widetilde x),
\qquad
\widetilde x(t_0)=\widetilde x_0.
\]
Тогда при той же константе $L$ для $f$
\begin{equation}
\label{eq:model-perturbation-bound-v01}
\begin{aligned}
\norm{x(t)-\widetilde x(t)}
&\leq
 e^{L(t-t_0)}\norm{x_0-\widetilde x_0}\\
&\quad+
\int_{t_0}^{t}
 e^{L(t-s)}
 \norm{f(s,\widetilde x(s))-\widetilde f(s,\widetilde x(s))}
\,\dd s.
\end{aligned}
\end{equation}
Первая часть --- усиленная ошибка начального состояния, вторая --- накопленная ошибка закона динамики. Формула отделяет два источника неточности и показывает, почему малой локальной ошибки модели недостаточно на длинной траектории.

\begin{application}[гибридная физико-нейросетевая модель]
В дифференцируемом симуляторе удобно разделить известную физику и обучаемую поправку:
\begin{equation}
\label{eq:hybrid-neural-ode-v01}
\dot x
=
f_{\mathrm{phys}}(t,x;\eta)
+r_\theta(t,x,\eta),
\qquad
x(t_0)=x_0(\eta).
\end{equation}
Здесь $\eta$ задаёт конструкцию, $f_{\mathrm{phys}}$ описывает известные законы, а $r_\theta$ компенсирует неучтённые эффекты по данным. Если обе части локально липшицевы по $x$, модель задаёт единственную траекторию. Оценка~\eqref{eq:model-perturbation-bound-v01} связывает ошибку обученной поправки с ошибкой конечного состояния, а зависимость $x_0(\eta)$ подчёркивает, что параметры проекта могут менять и начальную конфигурацию, и сам закон движения.
\end{application}

Для задач генеративного проектирования это даёт первый критерий пригодности динамической модели. До оптимизации и обратного распространения нужно проверить, что для допустимых параметров $\eta$ существует единственная траектория до момента оценки $T$, она остаётся в физически допустимой области, а малые ошибки данных и модели не усиливаются до неприемлемого уровня. Следующие разделы дадут более точные инструменты: спектр линейной системы, функции Ляпунова, устойчивость численной схемы и уравнения чувствительности.

\section{Линейная динамика и спектральные режимы}
\label{sec:linear-dynamics-spectral-modes}

\subsection{Почему линейная система является базовой моделью}

Линейная динамика возникает в двух разных ситуациях. Иногда закон процесса изначально линеен, например в малых тепловых, электрических или механических моделях. Чаще нелинейное уравнение рассматривают вблизи выбранного состояния и заменяют первым членом разложения Тейлора. Поэтому система
\begin{equation}
\label{eq:linear-autonomous-v02}
\dot x(t)=Ax(t),
\qquad
x(t_0)=x_0,
\end{equation}
где $A\in\R^{n\times n}$, одновременно служит точной моделью простого процесса и локальной моделью сложного процесса.

Матрица $A$ не переводит начальное состояние сразу в конечное. Она задаёт мгновенную скорость: в точке $x$ вектор $Ax$ показывает, куда и насколько быстро движется система. Конечный переход от $t_0$ к $t$ описывается другим оператором.

\begin{definition}
Матричной экспонентой называется сходящийся ряд
\begin{equation}
\label{eq:matrix-exponential-definition-v02}
\exp(A)
=
I+A+\frac{A^2}{2!}+\frac{A^3}{3!}+\cdots.
\end{equation}
Для $t\in\R$ используется обозначение $e^{tA}=\exp(tA)$.
\end{definition}

Почленное дифференцирование ряда даёт
\[
\frac{\dd}{\dd t}e^{tA}=Ae^{tA}=e^{tA}A,
\qquad
 e^{0A}=I.
\]
Следовательно, единственное решение задачи~\eqref{eq:linear-autonomous-v02} имеет вид
\begin{equation}
\label{eq:linear-autonomous-solution-v02}
x(t)=e^{(t-t_0)A}x_0.
\end{equation}
Матрица $\Phi(t,t_0)=e^{(t-t_0)A}$ называется оператором перехода состояния: она переносит любое начальное состояние в состояние в момент $t$.

Для одной и той же матрицы $A$ выполняются свойства
\begin{equation}
\label{eq:matrix-exponential-flow-properties-v02}
e^{(t+s)A}=e^{tA}e^{sA},
\qquad
\bigl(e^{tA}\bigr)^{-1}=e^{-tA}.
\end{equation}
Первое равенство выражает согласованность движения во времени: переход на $s$ единиц, а затем на $t$ единиц совпадает с одним переходом на $t+s$. Второе показывает, что точный линейный поток обратим для любого конечного времени, даже когда траектории быстро сжимаются.

\begin{caution}
В общем случае $e^{A+B}\neq e^Ae^B$. Равенство справедливо, если $AB=BA$. Поэтому разложение сложной динамики на последовательно применяемые части требует проверки коммутативности или внесения ошибки расщепления.
\end{caution}

При замене координат $x=Sy$, где $S$ обратима, матрица системы становится $S^{-1}AS$, а оператор перехода преобразуется согласованно:
\begin{equation}
\label{eq:matrix-exponential-coordinate-change-v02}
e^{tS^{-1}AS}=S^{-1}e^{tA}S.
\end{equation}
Физическое движение не зависит от выбранного базиса; меняются только координаты его описания.

\subsection{Собственные направления как независимые режимы}

Пусть $Av=\lambda v$. Если начальное состояние направлено вдоль $v$, то
\begin{equation}
\label{eq:eigenmode-evolution-v02}
x(t)=e^{\lambda(t-t_0)}v.
\end{equation}
Собственный вектор задаёт режим, форма которого не меняется, а собственное значение задаёт его временной закон.

Если $\lambda\in\R$, знак $\lambda$ различает затухание, сохранение и рост. Для комплексного собственного значения $\lambda=a+i\omega$ имеем
\[
e^{\lambda t}=e^{at}\bigl(\cos(\omega t)+i\sin(\omega t)\bigr).
\]
Вещественная часть $a$ определяет изменение амплитуды, а мнимая часть $\omega$ --- частоту вращения или колебания. У вещественной матрицы комплексные значения возникают сопряжёнными парами и вместе задают вещественное колебательное движение.

Если матрица диагонализируема,
\[
A=S\Lambda S^{-1},
\qquad
\Lambda=\operatorname{diag}(\lambda_1,\ldots,\lambda_n),
\]
то
\begin{equation}
\label{eq:diagonalizable-linear-flow-v02}
e^{tA}
=
S\operatorname{diag}\bigl(e^{t\lambda_1},\ldots,e^{t\lambda_n}\bigr)S^{-1}.
\end{equation}
Разложив $x_0=\sum_j c_jv_j$, получаем
\begin{equation}
\label{eq:modal-decomposition-v02}
x(t)=\sum_{j=1}^{n}c_j e^{\lambda_j(t-t_0)}v_j.
\end{equation}
Таким образом, линейная система раскладывается на режимы, которые развиваются независимо, а наблюдаемая траектория является их суммой.

\begin{definition}
Спектральной абсциссой матрицы $A$ называется число
\begin{equation}
\label{eq:spectral-abscissa-v02}
\alpha(A)=\max_{\lambda\in\sigma(A)}\Re\lambda.
\end{equation}
Оно показывает наибольшую асимптотическую скорость роста среди спектральных режимов.
\end{definition}

\begin{theorem}[спектральный критерий асимптотической устойчивости]
\label{thm:linear-spectral-stability-v02}
Нулевое состояние системы $\dot x=Ax$ асимптотически устойчиво тогда и только тогда, когда
\[
\Re\lambda<0
\qquad
\text{для всех }\lambda\in\sigma(A).
\]
В этом случае существуют $C>0$ и $\gamma>0$ такие, что
\begin{equation}
\label{eq:linear-exponential-decay-v02}
\norm{e^{tA}}\leq Ce^{-\gamma t},
\qquad t\geq0.
\end{equation}
\end{theorem}

Условие $\alpha(A)<0$ означает, что каждый режим затухает. Если хотя бы одно собственное значение имеет положительную вещественную часть, начальные данные с компонентой в соответствующем направлении растут. При нулевой вещественной части одной проверки знаков недостаточно: жорданов блок может породить полиномиальный рост.

Пространство состояний можно разложить на инвариантные части, порождённые обобщёнными собственными векторами с отрицательными, положительными и нулевыми вещественными частями собственных значений. Их называют устойчивым, неустойчивым и центральным подпространствами. Такое разложение показывает не только общий исход, но и какие компоненты начального состояния исчезают, растут или сохраняются.

\subsection{Жордановы блоки и временное усиление}

Если собственных векторов недостаточно для диагонализации, используется жорданов блок
\[
J=\lambda I+N,
\qquad
N^m=0.
\]
Поскольку $\lambda I$ и $N$ коммутируют,
\begin{equation}
\label{eq:jordan-exponential-v02}
e^{tJ}
=
e^{\lambda t}
\left(
I+tN+\frac{t^2N^2}{2!}+\cdots+
\frac{t^{m-1}N^{m-1}}{(m-1)!}
\right).
\end{equation}
Экспоненциальный множитель определяет далёкое будущее, но полиномиальный множитель может существенно влиять на конечном интервале.

Более того, даже диагонализируемая матрица с отрицательными вещественными частями собственных значений способна временно усиливать некоторые состояния, если её собственные направления почти линейно зависимы. Рассмотрим
\begin{equation}
\label{eq:nonnormal-transient-example-v02}
A=
\begin{pmatrix}
-1&\kappa\\
0&-2
\end{pmatrix},
\qquad
x(0)=
\begin{pmatrix}
0\\1
\end{pmatrix}.
\end{equation}
Оба собственных значения отрицательны, однако решение равно
\begin{equation}
\label{eq:nonnormal-transient-solution-v02}
x(t)=
\begin{pmatrix}
\kappa\bigl(e^{-t}-e^{-2t}\bigr)\\
e^{-2t}
\end{pmatrix}.
\end{equation}
Первая координата сначала возрастает и при $t=\log2$ достигает $\kappa/4$, после чего затухает. При большом $\kappa$ кратковременное усиление может быть большим, хотя система асимптотически устойчива.

\begin{mlconnection}
Для линейного непрерывного слоя отображение начального скрытого состояния имеет якобиан
\begin{equation}
\label{eq:linear-flow-input-jacobian-v02}
\frac{\partial x(T)}{\partial x_0}=e^{(T-t_0)A}.
\end{equation}
Поэтому норма $\norm{e^{(T-t_0)A}}$ одновременно измеряет усиление возмущений прямого прохода и возможный масштаб градиента по начальному состоянию. Знаки вещественных частей собственных значений описывают асимптотику, но для конечного горизонта важны также неортогональность собственных направлений и сингулярные числа оператора перехода. Именно здесь возникают аналоги затухающих и взрывающихся состояний в рекуррентных моделях.
\end{mlconnection}

\subsection{Внешнее воздействие и вариация постоянных}

Реальная система почти всегда получает нагрузку, управление или поток энергии. Для уравнения
\begin{equation}
\label{eq:forced-linear-system-v02}
\dot x(t)=Ax(t)+b(t),
\qquad
x(t_0)=x_0,
\end{equation}
решение задаётся формулой вариации постоянных
\begin{equation}
\label{eq:variation-of-constants-v02}
x(t)
=
e^{(t-t_0)A}x_0
+
\int_{t_0}^{t}e^{(t-s)A}b(s)\,\dd s.
\end{equation}
Первое слагаемое является свободным движением, второе суммирует отклики на воздействия, приложенные в разные моменты. Воздействие в момент $s$ сначала добавляет $b(s)\,\dd s$, а затем переносится оператором $e^{(t-s)A}$ до момента наблюдения.

Если $b(t)\equiv b$ и матрица $A$ обратима, равновесие удовлетворяет
\begin{equation}
\label{eq:linear-forced-equilibrium-v02}
Ax_*+b=0,
\qquad
x_*=-A^{-1}b.
\end{equation}
После замены $z=x-x_*$ получаем $\dot z=Az$. Поэтому спектр $A$ отвечает за устойчивость отклонений от вынужденного равновесия.

\begin{application}[тепловая модель проектируемой конструкции]
Пусть $\vartheta(t)\in\R^n$ содержит превышения температур в узлах конструкции над температурой среды. После пространственной дискретизации простая модель теплопереноса имеет вид
\begin{equation}
\label{eq:thermal-design-linear-model-v02}
C(\eta)\dot\vartheta
=
-K(\eta)\vartheta+Bq(t),
\end{equation}
где $C(\eta)$ --- матрица теплоёмкостей, $K(\eta)$ --- матрица теплопроводности, $q(t)$ --- тепловая нагрузка, а $\eta$ задаёт геометрию и материалы. Если $C$ положительно определена, то
\[
\dot\vartheta
=A(\eta)\vartheta+C(\eta)^{-1}Bq(t),
\qquad
A(\eta)=-C(\eta)^{-1}K(\eta).
\]
Собственные режимы $A(\eta)$ описывают независимые формы остывания. Собственное значение, ближайшее к нулю, задаёт самый медленный режим и часто определяет время выхода на допустимую температуру. При генеративном проектировании изменение каналов, толщин и материалов меняет $C$, $K$, спектр и тем самым весь переходный процесс, а не только стационарную температуру.
\end{application}

\subsection{От непрерывного потока к слою сети}

Если система наблюдается через равные интервалы длины $h$, точная динамика~\eqref{eq:linear-autonomous-v02} удовлетворяет
\begin{equation}
\label{eq:exact-linear-discrete-step-v02}
x_{k+1}=e^{hA}x_k.
\end{equation}
Поэтому линейная рекуррентная модель $x_{k+1}=Wx_k$ является дискретным аналогом линейного потока; при точном соответствии $W=e^{hA}$. Собственные значения перехода равны $e^{h\lambda_j}$, так что условие $\Re\lambda_j<0$ переходит в $|e^{h\lambda_j}|<1$.

Остаточный шаг
\begin{equation}
\label{eq:euler-residual-bridge-v02}
x_{k+1}=x_k+hAx_k=(I+hA)x_k
\end{equation}
заменяет точную экспоненту её первым приближением $e^{hA}=I+hA+O(h^2)$. Это даёт прямой мост между ОДУ и residual-архитектурой, но одновременно показывает различие: устойчивый непрерывный процесс может стать неустойчивым после слишком грубой дискретизации. Точная роль шага и численной схемы будет разобрана далее.

Линейная модель даёт четыре практических объекта, которые будут использоваться в остальных разделах: оператор перехода $e^{tA}$, спектральные режимы, оценку конечновременного усиления и свёртку внешнего воздействия с переходным оператором. В нелинейной системе те же объекты возникают локально после линеаризации, а в обучаемой динамике определяют распространение состояния и градиента.

\section{Нелинейная динамика, устойчивость и функции Ляпунова}
\label{sec:nonlinear-stability-lyapunov}

\subsection{Равновесие как рабочий режим}

Рассмотрим автономную систему
\begin{equation}
\label{eq:nonlinear-autonomous-v03}
\dot x=f(x;\eta),
\qquad
x\in\R^n,
\end{equation}
где параметр $\eta$ описывает конструкцию, материал или обучаемую модель. В отличие от линейного случая, скорость теперь меняется вместе с состоянием, поэтому одна матрица уже не описывает динамику во всём пространстве.

\begin{definition}
Точка $x_*=x_*(\eta)$ называется \emph{равновесием}, если
\begin{equation}
\label{eq:equilibrium-v03}
f(x_*;\eta)=0.
\end{equation}
Траектория, начатая в $x_*$, остаётся там: $x(t)\equiv x_*$.
\end{definition}

Равновесие не обязано быть нулём. В тепловой модели это стационарное распределение температуры, в механике --- положение покоя, в непрерывной нейронной модели --- скрытое состояние, которое слой больше не изменяет. После замены $z=x-x_*$ равновесие переносится в начало координат.

Само наличие равновесия не говорит, что система к нему возвращается после ошибки, возмущения или неточного начального состояния. Для этого нужно понятие устойчивости.

\begin{definition}
Равновесие $x_*$ системы~\eqref{eq:nonlinear-autonomous-v03} называется:

\emph{устойчивым по Ляпунову}, если для каждого $\varepsilon>0$ существует $\delta>0$ такое, что из $\norm{x(0)-x_*}<\delta$ следует
\[
\norm{x(t)-x_*}<\varepsilon
\qquad
\text{для всех }t\geq0;
\]

\emph{притягивающим}, если все достаточно близкие начальные состояния удовлетворяют $x(t)\to x_*$ при $t\to\infty$;

\emph{асимптотически устойчивым}, если оно одновременно устойчиво и притягивает;

\emph{экспоненциально устойчивым}, если существуют $C,\alpha,r>0$ такие, что
\begin{equation}
\label{eq:exponential-stability-v03}
\norm{x(t)-x_*}
\leq
C e^{-\alpha t}\norm{x(0)-x_*}
\end{equation}
для всех $\norm{x(0)-x_*}<r$ и $t\geq0$.
\end{definition}

Устойчивость означает отсутствие самопроизвольного усиления малой ошибки, притяжение --- возвращение к режиму, а экспоненциальная оценка дополнительно задаёт скорость возвращения. В инженерной задаче именно $\alpha^{-1}$ имеет смысл характерного времени затухания.

\subsection{Линеаризация около равновесия}

Пусть $f$ непрерывно дифференцируема по $x$. Запишем $x=x_*+z$. Формула Тейлора даёт
\begin{equation}
\label{eq:nonlinear-linearization-v03}
\dot z
=
Az+r(z),
\qquad
A=D_xf(x_*;\eta),
\qquad
\frac{\norm{r(z)}}{\norm{z}}\to0
\quad(z\to0).
\end{equation}
Матрица $A$ является якобианом векторного поля в рабочей точке. Она описывает первый порядок изменения скорости при малом отклонении от равновесия.

\begin{theorem}[устойчивость по линеаризации]
\label{thm:nonlinear-linearization-stability-v03}
Пусть $f\in C^1$ и $f(x_*)=0$.
Если все собственные значения $A=Df(x_*)$ имеют отрицательные вещественные части, то $x_*$ экспоненциально устойчиво.
Если хотя бы одно собственное значение имеет положительную вещественную часть, то $x_*$ неустойчиво.
\end{theorem}

Так локальный анализ нелинейной системы сводится к спектральному анализу из предыдущего раздела. Если спектр не пересекает мнимую ось, равновесие называется \emph{гиперболическим}; в достаточно малой окрестности линейная и нелинейная системы имеют одинаковую качественную картину устойчивых и неустойчивых направлений.

\begin{caution}[нулевая вещественная часть]
Если у $Df(x_*)$ есть собственное значение с нулевой вещественной частью, линеаризация может не дать ответа. Например,
\[
\dot x=-x^3
\]
имеет асимптотически устойчивое равновесие $0$, а
\[
\dot x=x^3
\]
--- неустойчивое. В обоих случаях $f'(0)=0$ и линейная система имеет вид $\dot z=0$.
\end{caution}

\begin{application}[устойчивость проектного режима]
Пусть для каждого проекта $\eta$ рабочее состояние определяется уравнением
\[
f(x_*(\eta);\eta)=0.
\]
Локальная устойчивость зависит от матрицы
\[
A(\eta)=D_xf(x_*(\eta);\eta).
\]
При изменении геометрии равновесие и спектр меняются одновременно. Поэтому проверка только номинального проекта недостаточна: полезно контролировать спектральный запас
\begin{equation}
\label{eq:stability-margin-design-v03}
\gamma(\eta)
=
-\max_{\lambda\in\sigma(A(\eta))}\Re\lambda.
\end{equation}
Условие $\gamma(\eta)>0$ означает локальную экспоненциальную устойчивость, а малое значение $\gamma$ предупреждает о медленном восстановлении и чувствительности к ошибкам модели.
\end{application}

\subsection{Функция Ляпунова: энергия без явного решения}

Линеаризация локальна и требует вычисления спектра. Другой подход состоит в поиске скалярной величины, которая убывает вдоль каждой траектории.

\begin{definition}
Функция $V$ называется \emph{функцией Ляпунова} для равновесия $x_*$ в окрестности $U$, если
\begin{equation}
\label{eq:lyapunov-positive-v03}
V(x_*)=0,
\qquad
V(x)>0
\quad
(x\in U\setminus\{x_*\}),
\end{equation}
а производная вдоль траекторий
\begin{equation}
\label{eq:lyapunov-orbital-derivative-v03}
\dot V(x)
=
\inner{\grad V(x)}{f(x)}
\end{equation}
неположительна.
\end{definition}

Здесь $V$ играет роль расстояния, энергии или штрафа за отклонение, но не обязана совпадать с евклидовой нормой. Важна не величина сама по себе, а знак её изменения по закону системы.

\begin{theorem}[прямой метод Ляпунова]
\label{thm:direct-lyapunov-v03}
Если в окрестности $x_*$ функция $V$ положительно определена и $\dot V\leq0$, то равновесие устойчиво. Если дополнительно $\dot V<0$ при $x\neq x_*$, то равновесие асимптотически устойчиво.
\end{theorem}

Идея проста. Малое подуровневое множество
\[
\Omega_c
=
\{x:V(x)\leq c\}
\]
окружает равновесие. Поскольку $V$ не возрастает, траектория, попавшая в $\Omega_c$, не может пересечь его границу наружу. При строгом убывании она вынуждена двигаться к точке, где убывание прекращается.

Если $\dot V\leq0$, но равенство возможно не только в $x_*$, применяется инвариантный принцип ЛаСалля: траектории стремятся к наибольшему инвариантному множеству внутри $\{x:\dot V(x)=0\}$. Для практики это означает, что нужно проверить не просто знак диссипации, а какие движения могут сохраняться при нулевой диссипации.

\subsection{Энергия затухающей механической системы}

Рассмотрим нелинейную конструкцию без внешней нагрузки:
\begin{equation}
\label{eq:nonlinear-mechanical-system-v03}
M(\eta)\ddot q
+C(\eta)\dot q
+\grad_q U(q;\eta)=0,
\end{equation}
где $M$ положительно определена, $C$ неотрицательно определена, а $U$ --- потенциальная энергия. Пусть $q_*$ является строгим локальным минимумом $U$, то есть $\grad U(q_*;\eta)=0$.

Для состояния $x=(q,v)$ естественная функция Ляпунова равна превышению полной энергии над энергией покоя:
\begin{equation}
\label{eq:mechanical-lyapunov-energy-v03}
V(q,v)
=
\frac12 v^\top Mv
+U(q;\eta)-U(q_*;\eta).
\end{equation}
Вблизи $(q_*,0)$ она положительна. Дифференцируя вдоль решения и используя~\eqref{eq:nonlinear-mechanical-system-v03}, получаем
\begin{align}
\dot V
&=
v^\top M\dot v
+\inner{\grad U(q;\eta)}{\dot q}
\notag\\
&=
v^\top\bigl(-Cv-\grad U(q;\eta)\bigr)
+\inner{\grad U(q;\eta)}{v}
\notag\\
&=
-v^\top Cv
\leq0.
\label{eq:mechanical-energy-dissipation-v03}
\end{align}

\begin{application}[что проверять при генеративном проектировании]
Формула~\eqref{eq:mechanical-energy-dissipation-v03} переводит устойчивость в свойства проектных матриц и энергии. Матрица массы должна оставаться положительно определённой, демпфирование --- неотрицательным, а желаемая конфигурация --- локальным минимумом потенциальной энергии. При положительно определённом $C$ энергия строго убывает вне покоя. При вырожденном демпфировании нужно дополнительно исключить недемпфированные движения с $v\neq0$.

Такая проверка сильнее теста одной траектории: она одновременно охватывает целое множество близких начальных состояний и возмущений.
\end{application}

\subsection{Квадратичные сертификаты и матричное уравнение Ляпунова}

Для линейной системы $\dot z=Az$ удобно искать
\begin{equation}
\label{eq:quadratic-lyapunov-v03}
V(z)=z^\top Pz,
\qquad
P=P^\top\succ0.
\end{equation}
Тогда
\begin{equation}
\label{eq:quadratic-lyapunov-derivative-v03}
\dot V
=
z^\top(A^\top P+PA)z.
\end{equation}
Если для некоторой $Q=Q^\top\succ0$ выполнено
\begin{equation}
\label{eq:continuous-lyapunov-equation-v03}
A^\top P+PA=-Q,
\end{equation}
то $\dot V=-z^\top Qz<0$ вне нуля.

\begin{theorem}[квадратичный критерий]
\label{thm:quadratic-lyapunov-criterion-v03}
Матрица $A$ имеет только собственные значения с отрицательными вещественными частями тогда и только тогда, когда для каждого $Q\succ0$ уравнение~\eqref{eq:continuous-lyapunov-equation-v03} имеет единственное решение $P\succ0$.
\end{theorem}

Матрица $P$ задаёт эллипсоидальные уровни энергии, приспособленные к динамике. Это особенно важно для неортогональных систем: евклидова норма может временно расти, хотя подходящая квадратичная энергия монотонно убывает.

\begin{mlconnection}[устойчивость обучаемого векторного поля]
Для Neural ODE или гибридной модели локальный якобиан
\[
J_\theta(x)=D_xf_\theta(x)
\]
управляет распространением малых возмущений. Штраф за положительную спектральную абсциссу может улучшить локальную устойчивость, но сам по себе не гарантирует её во всей рабочей области. Более сильный путь --- построить или обучить функцию $V_\phi(x)$ и проверять условия
\[
V_\phi(x)>0,
\qquad
\inner{\grad V_\phi(x)}{f_\theta(x)}<0
\]
на интересующем множестве состояний. Тогда нейронная сеть используется не только как аппроксиматор динамики, но и как кандидат на проверяемый сертификат устойчивости.
\end{mlconnection}

\subsection{Область притяжения и допустимые состояния}

Локальная асимптотическая устойчивость не означает, что к равновесию придёт любая траектория. \emph{Областью притяжения} называется множество начальных состояний, решения из которых существуют при всех $t\geq0$ и стремятся к $x_*$. Обычно найти её точно трудно, но функция Ляпунова даёт внутреннюю оценку.

Если подуровневое множество $\Omega_c=\{V\leq c\}$ компактно, содержит $x_*$ и на нём выполнено $\dot V<0$ вне равновесия, то $\Omega_c$ лежит в области притяжения. В задаче проектирования дополнительно требуется
\begin{equation}
\label{eq:safe-lyapunov-sublevel-v03}
\Omega_c
\subseteq
\mathcal X_{\mathrm{adm}},
\end{equation}
где $\mathcal X_{\mathrm{adm}}$ задаёт ограничения на напряжения, температуру, перемещения или другие физические величины.

Таким образом, устойчивость отвечает на вопрос, вернётся ли система к рабочему режиму, а включение~\eqref{eq:safe-lyapunov-sublevel-v03} --- останется ли переходный процесс допустимым. Для генеративного проектирования полезен не только лучший номинальный отклик, но и достаточно большая сертифицированная область безопасного восстановления.

Линеаризация и функции Ляпунова дополняют друг друга. Якобиан быстро описывает локальные направления и скорость затухания; функция Ляпунова даёт нелинейный энергетический сертификат и оценку области притяжения. Следующий раздел перейдёт от непрерывной модели к вычисляемой траектории и покажет, как выбор численной схемы может сохранить или разрушить эти свойства.

\section{Численное интегрирование и дискретная динамика}
\label{sec:numerical-integration-discrete-dynamics-v04}

Дифференциальное уравнение задаёт непрерывную траекторию, но в вычислении доступны лишь значения состояния в конечном наборе моментов времени. Поэтому практическая модель состоит не только из векторного поля $f$, но и из численной схемы, шага, критерия точности и правил остановки. Неудачная дискретизация способна исказить скорость затухания, создать ложную неустойчивость или скрыть быстрое движение, хотя исходная непрерывная система сформулирована корректно.

Пусть требуется решить задачу Коши
\begin{equation}
\label{eq:ivp-for-numerical-integration-v04}
\dot x(t)=f(t,x(t)),
\qquad
x(t_0)=x_0,
\qquad
x(t)\in\R^n.
\end{equation}
Выберем сетку
\[
t_k=t_0+kh,
\qquad
k=0,1,\ldots,N,
\qquad
h=\frac{T-t_0}{N}.
\]
Численный метод строит приближения $x_k\approx x(t_k)$ по рекуррентному правилу.

\subsection{Одношаговая схема и порядок точности}

\begin{definition}[одношаговый метод]
\label{def:one-step-method-v04}
Одношаговой схемой называется отображение $\Psi_h$, для которого
\begin{equation}
\label{eq:one-step-method-v04}
x_{k+1}=\Psi_h(t_k,x_k).
\end{equation}
Следующее состояние вычисляется только из текущего состояния, текущего времени и шага $h$.
\end{definition}

Точное решение также задаёт одношаговое отображение
\[
x(t_{k+1})=\Phi_h(t_k,x(t_k)),
\]
где $\Phi_h$ --- точный поток за время $h$. Численная схема заменяет недоступное $\Phi_h$ вычислимым $\Psi_h$.

\begin{definition}[локальная и глобальная ошибки]
\label{def:local-global-error-v04}
Локальная ошибка одного шага равна
\begin{equation}
\label{eq:local-defect-v04}
\delta_{k+1}
=
\Phi_h(t_k,x(t_k))
-
\Psi_h(t_k,x(t_k)).
\end{equation}
Здесь метод стартует из точного состояния, поэтому $\delta_{k+1}$ измеряет ошибку самой формулы шага. Глобальная ошибка
\begin{equation}
\label{eq:global-error-v04}
e_k=x(t_k)-x_k
\end{equation}
содержит также ошибки, накопленные на всех предыдущих шагах.
\end{definition}

Метод имеет порядок $p$, если при достаточно гладком решении локальная ошибка имеет порядок $\bigO(h^{p+1})$, а глобальная --- $\bigO(h^p)$ на фиксированном промежутке времени.

\begin{theorem}[накопление локальной ошибки]
\label{thm:local-to-global-error-v04}
Пусть схема согласована с точным потоком с локальной ошибкой
\[
\norm{\delta_{k+1}}
\leq Ch^{p+1},
\]
а её зависимость от состояния удовлетворяет оценке
\[
\norm{\Psi_h(t,x)-\Psi_h(t,y)}
\leq
(1+Lh)\norm{x-y}.
\]
Тогда на фиксированном промежутке $[t_0,T]$
\begin{equation}
\label{eq:global-order-bound-v04}
\max_{0\leq k\leq N}\norm{e_k}
\leq
\frac{C}{L}\bigl(e^{L(T-t_0)}-1\bigr)h^p
\end{equation}
при $L>0$; при $L=0$ правая часть заменяется на $C(T-t_0)h^p$.
\end{theorem}

Действительно, ошибка удовлетворяет рекурсии
\[
\norm{e_{k+1}}
\leq
(1+Lh)\norm{e_k}+Ch^{p+1}.
\]
После $N\asymp h^{-1}$ шагов локальная величина порядка $h^{p+1}$ превращается в глобальную величину порядка $h^p$. Поэтому малость ошибки одного шага ещё не означает, что вся траектория вычислена с той же степенью точности.

\subsection{Явный метод Эйлера}

Из интегральной формы задачи Коши следует
\[
x(t_{k+1})
=
x(t_k)+\int_{t_k}^{t_{k+1}}f(s,x(s))\,\dd s.
\]
Если на коротком промежутке заменить подынтегральную функцию её значением в левом конце, получаем
\begin{equation}
\label{eq:explicit-euler-v04}
x_{k+1}
=
x_k+h f(t_k,x_k).
\end{equation}
Это явный метод: правая часть полностью известна до вычисления $x_{k+1}$.

Разложение Тейлора даёт
\[
x(t_k+h)
=
x(t_k)+h f(t_k,x(t_k))+\bigO(h^2).
\]
Следовательно, локальная ошибка Эйлера равна $\bigO(h^2)$, а глобальная --- $\bigO(h)$. Уменьшение шага в два раза обычно уменьшает глобальную ошибку примерно в два раза, пока не начинают доминировать округление и неточность вычисления $f$.

\begin{example}[один шаг]
Для $\dot x=-2x$, $x(0)=1$ и $h=0.1$ метод Эйлера даёт
\[
x_1=1+0.1(-2)=0.8,
\]
тогда как точное значение равно $e^{-0.2}\approx0.819$. Формула правильно передаёт затухание, но систематически заменяет экспоненту линейным множителем.
\end{example}

\subsection{Устойчивость численной схемы}

Порядок точности описывает поведение при $h\to0$, но не отвечает, допустим ли выбранный конечный шаг. Для этого используется тестовое уравнение
\begin{equation}
\label{eq:linear-test-equation-v04}
\dot x=\lambda x,
\qquad
\operatorname{Re}\lambda<0.
\end{equation}
Его точное решение за один шаг умножается на $e^{h\lambda}$ и затухает. Численная схема имеет вид
\begin{equation}
\label{eq:stability-function-v04}
x_{k+1}=R(h\lambda)x_k,
\end{equation}
где $R$ называется функцией устойчивости.

\begin{definition}[область абсолютной устойчивости]
\label{def:absolute-stability-region-v04}
Областью абсолютной устойчивости метода называется множество
\[
\mathcal S
=
\{z\in\C:\abs{R(z)}<1\}.
\]
Если $h\lambda\in\mathcal S$, численная последовательность затухает так же, как точное решение тестового уравнения.
\end{definition}

Для явного Эйлера
\[
R(z)=1+z.
\]
При вещественном $\lambda<0$ условие $\abs{1+h\lambda}<1$ равносильно
\begin{equation}
\label{eq:euler-step-stability-v04}
0<h<\frac{2}{\abs{\lambda}}.
\end{equation}
Если шаг больше, устойчивая непрерывная система превращается в осциллирующую или растущую дискретную последовательность.

\begin{caution}[точность и устойчивость --- разные требования]
Малый локальный дефект не спасает схему, если $h\lambda$ лежит вне области устойчивости. И наоборот, устойчивый шаг может быть слишком грубым для нужной точности. На практике должны одновременно выполняться оба требования.
\end{caution}

\subsection{Жёсткие системы}

Рассмотрим систему
\begin{equation}
\label{eq:stiff-two-scale-system-v04}
\dot x_1=-x_1,
\qquad
\dot x_2=-100x_2.
\end{equation}
Вторая компонента исчезает примерно за время $10^{-2}$, а первая меняется на масштабе порядка единицы. После короткого переходного процесса интерес представляет главным образом медленная компонента, однако явный Эйлер всё равно требует $h<0.02$ из-за быстрого режима.

\begin{definition}[жёсткость в вычислительном смысле]
\label{def:stiffness-practical-v04}
Систему называют жёсткой, когда шаг явного метода ограничивается быстрыми устойчивыми режимами значительно сильнее, чем того требует точность медленной интересующей динамики.
\end{definition}

Жёсткость не является просто наличием большого коэффициента. Важна комбинация сильно различающихся временных масштабов, длительного интервала интегрирования и ограниченной области устойчивости выбранной схемы. Такие ситуации возникают в теплопереносе, химической кинетике, контактных моделях, релаксации материалов и гибридных физических симуляторах.

\subsection{Метод Хойна и классическая схема Рунге--Кутты}

Метод Хойна сначала делает прогноз Эйлера
\begin{equation}
\label{eq:heun-predictor-v04}
y_k=x_k+h f(t_k,x_k),
\end{equation}
а затем усредняет наклоны в начале и в предсказанном конце шага:
\begin{equation}
\label{eq:heun-corrector-v04}
x_{k+1}
=
x_k+\frac h2
\bigl(
 f(t_k,x_k)+f(t_{k+1},y_k)
\bigr).
\end{equation}
Локальная ошибка имеет порядок $\bigO(h^3)$, глобальная --- $\bigO(h^2)$. За два вычисления векторного поля метод учитывает изменение направления внутри шага.

Классическая схема Рунге--Кутты четвёртого порядка использует четыре наклона:
\begin{align}
 k_1&=f(t_k,x_k),
 &
 k_2&=f\left(t_k+\frac h2,x_k+\frac h2k_1\right),
 \notag\\
 k_3&=f\left(t_k+\frac h2,x_k+\frac h2k_2\right),
 &
 k_4&=f(t_k+h,x_k+hk_3),
\label{eq:rk4-stages-v04}
\end{align}
после чего
\begin{equation}
\label{eq:rk4-update-v04}
x_{k+1}
=
x_k+\frac h6(k_1+2k_2+2k_3+k_4).
\end{equation}
При достаточной гладкости глобальная ошибка равна $\bigO(h^4)$. Высший порядок не означает автоматического превосходства: один шаг дороже, область устойчивости остаётся ограниченной, а на негладкой или шумной модели формальный порядок может не реализоваться.

\begin{computationalnote}[проверка порядка]
Если точное решение неизвестно, полезно вычислить траектории с шагами $h$, $h/2$ и $h/4$ и сравнить их на общих узлах. Для метода порядка $p$ отношение ошибок в асимптотическом режиме должно быть близко к $2^p$. Отсутствие такого поведения указывает на слишком крупный шаг, негладкость, жёсткость или ошибку реализации.
\end{computationalnote}

\subsection{Адаптивный шаг}

Один фиксированный шаг редко оптимален на всей траектории. В спокойной области его можно увеличить, а около быстрого перехода, контакта или резкого изменения нагрузки --- уменьшить. Простейшая оценка строится сравнением одного шага длины $h$ и двух шагов длины $h/2$.

Пусть $x_h$ и $x_{h/2}$ обозначают два приближения в одном и том же конечном моменте. Для метода порядка $p$ оценка ошибки более точного решения имеет вид
\begin{equation}
\label{eq:step-doubling-error-v04}
\widehat e
\approx
\frac{\norm{x_{h/2}-x_h}}{2^p-1}.
\end{equation}
Если $\widehat e$ превышает допуск, шаг отклоняется и повторяется с меньшим $h$. Если ошибка намного меньше допуска, следующий шаг можно увеличить.

\begin{caution}[малый шаг не всегда означает численную проблему]
При приближении к настоящему разрушению решения за конечное время адаптивный алгоритм также будет уменьшать шаг. Поэтому сообщение о слишком малом шаге может указывать как на жёсткость или неудачную схему, так и на реальную сингулярность модели. Различить эти случаи можно только с помощью анализа уравнения, контроля физических величин и сравнения нескольких методов.
\end{caution}

\subsection{Неявный Эйлер}

Неявная схема вычисляет новое состояние из уравнения
\begin{equation}
\label{eq:implicit-euler-v04}
x_{k+1}
=
x_k+h f(t_{k+1},x_{k+1}).
\end{equation}
Здесь $x_{k+1}$ присутствует по обе стороны, поэтому на каждом шаге требуется решить нелинейную систему. Для тестового уравнения~\eqref{eq:linear-test-equation-v04}
\[
R(z)=\frac{1}{1-z}.
\]
При любом $\operatorname{Re}z<0$ выполняется $\abs{R(z)}<1$. Значит, неявный Эйлер не теряет устойчивость на отрицательной вещественной полуоси при увеличении шага.

Цена этой устойчивости --- решение уравнения
\[
G(y)=y-x_k-hf(t_{k+1},y)=0.
\]
Метод Ньютона использует матрицу
\begin{equation}
\label{eq:implicit-newton-jacobian-v04}
D_yG(y)=I-hD_xf(t_{k+1},y).
\end{equation}
Таким образом, один неявный шаг может требовать нескольких вычислений $f$, якобиана и решения линейных систем. Для жёсткой задачи это часто выгоднее тысяч устойчивых явных шагов.

\begin{application}[тепловая модель конструкции]
После пространственной дискретизации теплопроводности часто возникает система
\[
C(\eta)\dot T+K(\eta)T=r(t),
\]
где $C$ --- матрица теплоёмкости, $K$ --- матрица теплопроводности, а $\eta$ задаёт геометрию и материал. Тонкие элементы и мелкая сетка создают быстро затухающие температурные моды. Явный шаг ограничивается наибольшим собственным значением $C^{-1}K$, тогда как неявный шаг сохраняет устойчивость и позволяет идти по медленному масштабу нагрева. При генеративном проектировании это различие влияет на стоимость оценки каждой кандидатной конструкции.
\end{application}

\subsection{Дискретизация может разрушить физическое свойство}

Рассмотрим гармонический осциллятор
\begin{equation}
\label{eq:harmonic-oscillator-v04}
\dot q=v,
\qquad
\dot v=-\omega^2 q.
\end{equation}
Его энергия
\[
E(q,v)=\frac12v^2+\frac12\omega^2q^2
\]
сохраняется. Явный Эйлер даёт
\begin{equation}
\label{eq:euler-harmonic-update-v04}
\begin{pmatrix}
q_{k+1}\\
v_{k+1}
\end{pmatrix}
=
\begin{pmatrix}
1 & h\\
-h\omega^2 & 1
\end{pmatrix}
\begin{pmatrix}
q_k\\v_k
\end{pmatrix}.
\end{equation}
Собственные значения матрицы шага равны $1\pm ih\omega$, а их модули равны
\[
\sqrt{1+h^2\omega^2}>1.
\]
Следовательно, численная амплитуда растёт при любом $h>0$, хотя точная траектория ограничена.

\begin{principle}[согласование схемы со структурой]
\label{pr:structure-preserving-discretization-v04}
При длительном моделировании важно проверять не только норму ошибки, но и свойства, которые схема должна сохранять или правильно рассеивать: энергию, массу, положительность, ограничения и монотонность функции Ляпунова. Универсально лучшего интегратора нет; выбор зависит от структуры задачи.
\end{principle}

Для консервативной механики применяются симплектические методы, для диссипативной системы --- схемы с корректным убыванием энергии, для положительных концентраций --- методы, не создающие отрицательные значения. Подробная теория таких методов выходит за пределы главы, но сам критерий выбора необходим для надёжного симулятора.

\subsection{Связь с остаточными сетями и Neural ODE}

Остаточный блок имеет форму
\begin{equation}
\label{eq:resnet-euler-v04}
x_{k+1}=x_k+h f_{\theta_k}(x_k).
\end{equation}
Это ровно шаг явного Эйлера для обучаемого векторного поля. При одинаковых параметрах $\theta_k=\theta$ все слои дискретизируют автономную систему; при разных параметрах получается неавтономная динамика, где глубина играет роль времени.

\begin{mlconnection}[дискретизация является частью модели]
Увеличение глубины при фиксированном горизонте $T$ уменьшает $h=T/N$ и меняет не только число параметров, но и качество приближения непрерывного потока. Большой residual-шаг может сделать устойчивое векторное поле неустойчивой сетью по той же причине, что и в условии~\eqref{eq:euler-step-stability-v04}. В Neural ODE выбор решателя, допусков и стратегии шага определяет фактическое отображение входа в выход и стоимость обучения. Поэтому их нельзя рассматривать как нейтральную техническую настройку.
\end{mlconnection}

При обучении полезно отделять три источника расхождения: ошибку модели $f_\theta$, ошибку численного интегрирования и ошибку оптимизации параметров. Уменьшение допуска решателя влияет только на вторую составляющую и не исправляет неверное векторное поле. С другой стороны, грубая схема способна замаскировать улучшение модели или породить ложный градиент.

\subsection{Практический протокол проверки}

Надёжная вычислительная модель проверяется последовательным уточнением шага и сравнением хотя бы двух схем разного порядка. На контрольных задачах с известным решением измеряется фактический порядок сходимости. На физической задаче дополнительно отслеживаются энергия, масса, положительность, ограничения и остаток уравнения. Для жёсткой системы сравниваются явный и неявный методы по точности на единицу вычислительной стоимости. Адаптивный решатель тестируется около быстрых переходов и на траекториях, близких к разрушению решения.

В генеративном проектировании этот протокол должен применяться не только к одной номинальной геометрии. Изменение конструкции меняет матрицы, спектр и временные масштабы, поэтому шаг, устойчивый для одного кандидата, может быть непригоден для другого. Следовательно, численная устойчивость и стоимость симуляции сами становятся свойствами проектного пространства.

Численная схема превращает непрерывный поток в конечную вычислительную композицию. Следующий раздел использует эту композицию для вычисления производных по параметрам и показывает разницу между прямой чувствительностью, обратным распространением через шаги и сопряжённым методом.

\section{Чувствительность решения и сопряжённый градиент}
\label{sec:sensitivity-and-adjoint-gradient}

В генеративном проектировании параметры $\theta\in\R^p$ задают геометрию, материал, закон управления или коэффициенты нейронной части модели. После выбора $\theta$ решается параметрическая задача Коши
\begin{equation}
\label{eq:parametric-ivp-v05}
\dot x(t)=f\bigl(t,x(t),\theta\bigr),
\qquad
x(t_0)=x_0(\theta),
\end{equation}
а затем по траектории вычисляется целевая функция. Поэтому оптимизатору требуется не только значение $J(\theta)$, но и производная $\nabla_\theta J$.

Здесь возникает основная вычислительная цепочка дифференцируемого симулятора:
\[
\theta
\longmapsto
x(\cdot;\theta)
\longmapsto
J(\theta)
\longmapsto
\nabla_\theta J.
\]
Вопрос состоит не в том, существует ли формальная производная, а в том, как вычислить её без построения огромного якобиана всей траектории.

\subsection{От производной решения к вариационному уравнению}

Предположим, что $f$ непрерывна по времени и непрерывно дифференцируема по $(x,\theta)$ в рабочей области. Тогда при стандартных условиях корректности решение локально гладко зависит от параметров. Для каждой координаты $\theta_j$ определим вектор чувствительности
\[
s_j(t)
=
\frac{\partial x(t;\theta)}{\partial\theta_j}
\in\R^n.
\]
Соберём эти векторы в матрицу
\begin{equation}
\label{eq:sensitivity-matrix-v05}
S(t)
=
\frac{\partial x(t;\theta)}{\partial\theta}
\in\R^{n\times p}.
\end{equation}

Дифференцируя~\eqref{eq:parametric-ivp-v05} по $\theta$ и применяя правило цепочки, получаем
\begin{equation}
\label{eq:forward-sensitivity-v05}
\dot S(t)
=
D_xf\bigl(t,x(t),\theta\bigr)S(t)
+
D_\theta f\bigl(t,x(t),\theta\bigr),
\qquad
S(t_0)=D_\theta x_0(\theta).
\end{equation}
Это \emph{прямое вариационное уравнение}. Оно линейно по неизвестной матрице $S$, хотя исходная динамика может быть нелинейной.

Формула~\eqref{eq:forward-sensitivity-v05} имеет прозрачный смысл. Матрица $D_xf$ переносит уже возникшее возмущение состояния вдоль траектории, а $D_\theta f$ непрерывно создаёт новое возмущение из-за изменения параметров. Начальное условие учитывает зависимость стартового состояния от $\theta$.

\begin{example}[однопараметрическая релаксация]
Рассмотрим
\[
\dot x=-\theta x,
\qquad
x(0)=x_0,
\qquad
\theta>0.
\]
Точное решение равно $x(t)=x_0e^{-\theta t}$. Чувствительность
\[
s(t)=\frac{\partial x}{\partial\theta}=-tx_0e^{-\theta t}
\]
должна удовлетворять
\[
\dot s=-\theta s-x,
\qquad
s(0)=0.
\]
Подстановка подтверждает равенство. Даже в этом простом примере чувствительность не совпадает с состоянием: она описывает, насколько изменится вся траектория при малом изменении коэффициента затухания.
\end{example}

\subsection{Градиент через прямую чувствительность}

Пусть качество конструкции измеряется функционалом
\begin{equation}
\label{eq:trajectory-objective-v05}
J(\theta)
=
\Phi\bigl(x(T;\theta),\theta\bigr)
+
\int_{t_0}^{T}
\ell\bigl(t,x(t;\theta),\theta\bigr)\,\dd t.
\end{equation}
Термин $\Phi$ оценивает конечное состояние, а $\ell$ --- накопленную стоимость: энергию, массу, перегрев, вибрацию или отклонение от желаемой траектории.

Правило цепочки даёт
\begin{equation}
\label{eq:direct-gradient-v05}
\nabla_\theta J
=
\nabla_\theta\Phi
+S(T)^\top\nabla_x\Phi
+
\int_{t_0}^{T}
\left(
\nabla_\theta\ell
+S(t)^\top\nabla_x\ell
\right)\dd t.
\end{equation}
Все производные в правой части вычисляются вдоль найденной траектории.

Прямой метод прост и надёжен: одновременно интегрируются $x$ и $S$, после чего применяется~\eqref{eq:direct-gradient-v05}. Однако размер $S$ равен $n\times p$. При тысячах или миллионах параметров хранение и перенос этой матрицы становятся дороже самой модели.

\begin{principle}[когда прямой метод уместен]
\label{pr:when-forward-sensitivity-v05}
Прямая чувствительность особенно удобна, когда параметров немного, а интересующих выходов много. Один прогон сразу даёт производные всей траектории по каждому параметру.
\end{principle}

\subsection{Зачем нужен сопряжённый метод}

Во многих задачах машинного обучения $p$ велико, но оптимизируется одна скалярная функция $J$. Вместо полного якобиана $S$ достаточно вычислить его произведение с градиентом потерь. Эту задачу решает сопряжённая переменная
\[
a(t)
=
\frac{\partial J}{\partial x(t)}
\in\R^n,
\]
которая показывает, как малое возмущение состояния в момент $t$ влияет на итоговую стоимость.

Для функционала~\eqref{eq:trajectory-objective-v05} сопряжённая система имеет вид
\begin{equation}
\label{eq:adjoint-ode-v05}
\dot a(t)
=
- D_xf\bigl(t,x(t),\theta\bigr)^\top a(t)
-
\nabla_x\ell\bigl(t,x(t),\theta\bigr),
\qquad
 a(T)=\nabla_x\Phi\bigl(x(T),\theta\bigr).
\end{equation}
Она интегрируется \emph{назад} от $T$ к $t_0$.

\begin{derivation}
Рассмотрим возмущение параметров $\delta\theta$ и соответствующее возмущение состояния $\delta x$. Линеаризация динамики даёт
\[
\delta\dot x
=D_xf\,\delta x+D_\theta f\,\delta\theta.
\]
Добавим к вариации цели тождественно нулевой интеграл
\[
\int_{t_0}^{T}
a(t)^\top
\bigl(D_xf\,\delta x+D_\theta f\,\delta\theta-\delta\dot x\bigr)
\,\dd t.
\]
Интегрирование члена $-a^\top\delta\dot x$ по частям переносит производную с $\delta x$ на $a$. Если выбрать конечное условие и динамику~\eqref{eq:adjoint-ode-v05}, все внутренние члены при $\delta x$ сокращаются. Остаются только члены при $\delta\theta$ и возмущение начального состояния.
\end{derivation}

В результате
\begin{equation}
\label{eq:adjoint-gradient-v05}
\nabla_\theta J
=
\nabla_\theta\Phi
+
D_\theta x_0(\theta)^\top a(t_0)
+
\int_{t_0}^{T}
\left(
\nabla_\theta\ell
+D_\theta f^\top a
\right)\dd t.
\end{equation}
Главное преимущество состоит в том, что размер $a(t)$ равен $n$ и не зависит от числа параметров $p$. Параметры входят только через векторно-якобианское произведение $D_\theta f^\top a$, которое автоматическое дифференцирование вычисляет без явного построения полного якобиана.

\begin{principle}[когда сопряжённый метод уместен]
\label{pr:when-adjoint-v05}
Сопряжённый метод особенно выгоден для одной или нескольких скалярных целей при большом числе параметров. Стоимость обратного прохода определяется главным образом размерностью состояния и числом вычислений векторного поля, а не числом столбцов полного якобиана чувствительности.
\end{principle}

\subsection{Конечная потеря и формула Chen et al.}

В базовой постановке Neural ODE потеря зависит только от конечного состояния:
\[
L=L\bigl(z(t_1)\bigr),
\qquad
\dot z=f(z,t,\theta).
\]
Тогда $\ell=0$, и сопряжённая система упрощается до
\begin{equation}
\label{eq:chen-adjoint-v05}
\dot a(t)
=-D_zf\bigl(z(t),t,\theta\bigr)^\top a(t),
\qquad
 a(t_1)=\nabla_zL\bigl(z(t_1)\bigr).
\end{equation}
Градиент по параметрам равен
\begin{equation}
\label{eq:chen-parameter-gradient-v05}
\nabla_\theta L
=
\int_{t_0}^{t_1}
D_\theta f\bigl(z(t),t,\theta\bigr)^\top a(t)\,\dd t.
\end{equation}
Именно эту схему Chen et al. используют для обратного распространения через чёрный ящик ODE-решателя: состояние, сопряжённая переменная и накопитель градиента объединяются в одну расширенную систему и интегрируются назад.

При обратном времени удобно ввести $g(t_1)=0$ и записать
\begin{equation}
\label{eq:augmented-adjoint-system-v05}
\frac{\dd}{\dd t}
\begin{pmatrix}
z\\a\\g
\end{pmatrix}
=
\begin{pmatrix}
f(z,t,\theta)\\
-D_zf(z,t,\theta)^\top a\\
-D_\theta f(z,t,\theta)^\top a
\end{pmatrix}.
\end{equation}
После интегрирования от $t_1$ к $t_0$ величина $g(t_0)$ совпадает с интегралом~\eqref{eq:chen-parameter-gradient-v05}.

Если потеря зависит от состояний в нескольких моментах $t_i$, сопряжённая переменная между наблюдениями решает то же ОДУ, а в каждом $t_i$ получает скачок
\begin{equation}
\label{eq:adjoint-jump-v05}
 a(t_i^-)
=
 a(t_i^+)
+
\nabla_x L_i\bigl(x(t_i)\bigr).
\end{equation}
Это непрерывный аналог обратного распространения через последовательность промежуточных выходов.

\subsection{Непрерывный и дискретный градиенты}

Есть два разных объекта, которые часто называют «градиентом через ODE». Первый получается, если сначала дискретизировать решатель и затем применить reverse-mode autodiff ко всем его операциям. Второй выводится для непрерывной задачи и затем численно интегрируется как сопряжённое ОДУ.

При конечном шаге и ненулевом допуске эти два градиента не обязаны совпадать точно. Дискретный градиент является производной конкретного вычислительного алгоритма, а непрерывный --- приближением производной точного потока. Разница уменьшается при согласованном уточнении решателей, но может быть заметна на жёстких, неустойчивых или плохо обратимых траекториях.

\begin{caution}[память обменивается на повторные вычисления]
Заявление о постоянной памяти в классической Neural ODE относится к схеме, где промежуточные состояния прямого прохода не сохраняются, а траектория восстанавливается обратным интегрированием. Это экономит память, но требует дополнительных вычислений и может накапливать ошибку восстановления. Практический компромисс --- хранить контрольные точки и переинтегрировать только между ними либо дифференцировать непосредственно через шаги решателя.
\end{caution}

Выбор метода определяется не лозунгом, а задачей. Для небольшой физической модели с десятком параметров прямые чувствительности проще проверять. Для нейронного векторного поля с миллионами весов нужен reverse-mode подход. Для жёсткой модели необходимо согласовать прямой и обратный решатели и отдельно проверить градиент конечными разностями на малой тестовой задаче.

\subsection{Проверка градиента}

Пусть $v\in\R^p$ --- случайное направление единичной нормы. Направленная производная должна удовлетворять
\begin{equation}
\label{eq:directional-gradient-check-v05}
\nabla_\theta J(\theta)^\top v
\approx
\frac{J(\theta+\varepsilon v)-J(\theta-\varepsilon v)}{2\varepsilon}.
\end{equation}
Слишком большое $\varepsilon$ даёт ошибку усечения, слишком малое --- потерю точности из-за округления и шума решателя. На практике проверяют несколько значений $\varepsilon$ и ищут диапазон, где расхождение убывает, а затем перестаёт улучшаться.

Эта проверка сравнивает полный вычислительный конвейер, включая ODE-решатель, целевую функцию и реализацию сопряжённого прохода. Она не доказывает правильность на всех параметрах, но быстро обнаруживает неверный знак, транспонирование, пропущенный вклад начального состояния или несовместимые допуски прямого и обратного интегрирования.

\subsection{Пример: динамическая оптимизация конструкции}

Пусть $q(t)$ --- перемещения конструкции, а параметры $\theta$ определяют толщины элементов и коэффициенты материала. Динамика имеет вид
\[
M(\theta)\ddot q+C(\theta)\dot q+K(\theta)q=p(t).
\]
После сведения к первому порядку получаем~\eqref{eq:parametric-ivp-v05}. Выберем гладкую цель
\begin{equation}
\label{eq:structural-dynamic-objective-v05}
J(\theta)
=
\frac12\norm{Wq(T)-q_{\mathrm{tar}}}^2
+
\frac\beta2\int_0^T\norm{Wq(t)}^2\,\dd t
+
\gamma m(\theta),
\end{equation}
где первый член контролирует конечное положение, второй подавляет вибрации, а $m(\theta)$ штрафует массу.

Прямой метод интегрирует производные $q$ и $\dot q$ по каждой толщине. Сопряжённый метод делает один прямой прогон состояния и один обратный прогон переменной $a$, после чего собирает вклады производных $M$, $C$, $K$ и $m$. Если параметров тысячи, разница в размере задачи становится принципиальной.

\begin{application}[градиент для генератора конструкций]
Пусть генератор с параметрами $\psi$ по латентному коду $\xi$ выдаёт проект
\[
\theta=G_\psi(\xi).
\]
Дифференцируемый симулятор возвращает $J(\theta)$ и $\nabla_\theta J$. Тогда правило цепочки даёт
\begin{equation}
\label{eq:generator-through-simulator-v05}
\nabla_\psi J
=
D_\psi G_\psi(\xi)^\top
\nabla_\theta J.
\end{equation}
Сопряжённый метод превращает физическую траекторию в градиентный сигнал для генератора. Так дифференциальное уравнение становится не только проверкой готовой конструкции, но и обучаемым звеном процесса проектирования.
\end{application}

\subsection{Итоговая вычислительная схема}

Для одной оценки параметров сначала решается прямая задача~\eqref{eq:parametric-ivp-v05} и вычисляется $J$. Затем выбирается способ дифференцирования. Прямая чувствительность решает~\eqref{eq:forward-sensitivity-v05}; сопряжённый метод решает~\eqref{eq:adjoint-ode-v05} назад и использует~\eqref{eq:adjoint-gradient-v05}. Полученный градиент проверяется направленной разностью~\eqref{eq:directional-gradient-check-v05} и только после этого передаётся оптимизатору.

Таким образом, чувствительность связывает локальные производные векторного поля с производной всей траектории, а сопряжённая переменная сжимает эту информацию до направления, необходимого скалярной цели. Следующий раздел использует этот аппарат для непрерывных нейронных слоёв, моделей временных рядов и генеративных потоков.

\section{Neural ODE и непрерывные генеративные модели}
\label{sec:neural-ode-and-continuous-generative-models}

\subsection{Непрерывная глубина как предел residual-сети}

Residual-блок преобразует скрытое состояние по правилу
\begin{equation}
\label{eq:residual-step-v06}
h_{k+1}
=
h_k+\Delta t\,f_\theta(h_k,t_k).
\end{equation}
Если интерпретировать номер слоя как дискретное время, то~\eqref{eq:residual-step-v06} является шагом явного Эйлера для уравнения
\begin{equation}
\label{eq:neural-ode-v06}
\dot h(t)=f_\theta\bigl(h(t),t\bigr),
\qquad
h(t_0)=h_0.
\end{equation}
Вместо заранее заданной последовательности слоёв модель задаёт векторное поле, а выход вычисляется как
\begin{equation}
\label{eq:odesolve-layer-v06}
h(t_1)
=
\operatorname{ODESolve}
\bigl(h_0,f_\theta,t_0,t_1;\tau\bigr),
\end{equation}
где $\tau$ обозначает настройки численного решателя: метод, абсолютный и относительный допуски, ограничения на шаг и число вычислений правой части.

\begin{definition}
\emph{Neural ODE-слоем} называется отображение начального состояния $h_0$ в конечное состояние $h(t_1)$, определённое задачей Коши~\eqref{eq:neural-ode-v06}, в которой векторное поле $f_\theta$ параметризовано нейронной сетью.
\end{definition}

Это определение разделяет три объекта. Нейронная сеть задаёт локальную скорость $f_\theta$; дифференциальное уравнение задаёт непрерывную траекторию; численный решатель приближённо вычисляет конечное отображение. Поэтому архитектура Neural ODE определяется не только весами $\theta$, но и классом допустимых траекторий и вычислительной политикой $\tau$.

В обычной глубокой сети число преобразований фиксировано архитектурой. Адаптивный решатель выбирает точки вычисления $f_\theta$ по ходу интегрирования. Сложная траектория требует больше обращений к сети, простая --- меньше. Число таких обращений обычно обозначают $\mathrm{NFE}$, number of function evaluations. Оно играет роль фактической глубины и одновременно определяет значительную часть стоимости прямого и обратного проходов.

\begin{principle}[допуск решателя является частью модели]
\label{pr:solver-tolerance-part-of-model-v06}
Изменение допуска меняет не только время вычисления, но и приближённое отображение $h_0\mapsto h(t_1)$. Поэтому точность интегрирования следует рассматривать как гиперпараметр модели и проверять отдельно при обучении и применении.
\end{principle}

Более строгий допуск обычно увеличивает $\mathrm{NFE}$ и уменьшает численную ошибку. Слишком грубый допуск ускоряет вычисление, но может заметно изменить признаки и градиенты. Слишком строгий допуск не гарантирует лучшего качества данных, зато способен сделать обучение неоправданно дорогим.

\subsection{Что даёт корректность задачи Коши}

Если $f_\theta$ локально липшицево по состоянию и решение существует на всём $[t_0,t_1]$, то каждому $h_0$ соответствует единственная траектория. Поток
\[
\Phi_{t_1,t_0}:h_0\longmapsto h(t_1)
\]
обратим: обратное отображение получается интегрированием той же динамики от $t_1$ к $t_0$. Это не дополнительная эвристика Neural ODE, а прямое следствие единственности решения.

Обратимость полезна, когда требуется сохранять информацию или строить поток вероятности. Одновременно она ограничивает геометрию преобразования: две различные траектории не могут встретиться в один момент и продолжиться как одна траектория. Если задача требует необратимого сжатия информации, его приходится реализовывать кодировщиком, декодировщиком, изменением размерности или отдельным дискретным слоем.

\begin{caution}[не всякая глубокая сеть имеет удобный непрерывный аналог]
Непрерывный поток естественен, когда данные или скрытая модель действительно описывают эволюцию. Если задача состоит из нескольких качественно разных операций, фиксированная композиция дискретных блоков может быть проще, дешевле и выразительнее. Neural ODE --- не замена всем residual-сетям, а архитектура с конкретным структурным предположением о плавной динамике.
\end{caution}

Липшицева константа векторного поля влияет сразу на несколько свойств. Большие значения усиливают чувствительность траектории к возмущениям, затрудняют численное интегрирование и могут увеличивать $\mathrm{NFE}$. Поэтому ограничения нормы весов, спектральная нормализация и штрафы на якобиан имеют здесь не только статистический, но и динамический смысл: они контролируют деформацию потока.

\subsection{Обучение непрерывного слоя}

Пусть модель вычисляет
\[
h_1=\Phi_\theta(h_0),
\qquad
L=L(h_1).
\]
Прямой проход решает~\eqref{eq:neural-ode-v06}. Обратный проход использует сопряжённую переменную из предыдущего раздела:
\begin{equation}
\label{eq:neural-ode-adjoint-v06}
\dot a(t)
=
-D_hf_\theta\bigl(h(t),t\bigr)^\top a(t),
\qquad
 a(t_1)=\nabla_hL\bigl(h(t_1)\bigr),
\end{equation}
и вычисляет
\begin{equation}
\label{eq:neural-ode-gradient-v06}
\nabla_\theta L
=
\int_{t_0}^{t_1}
D_\theta f_\theta\bigl(h(t),t\bigr)^\top a(t)\,\dd t.
\end{equation}
Chen et al. объединяют состояние, сопряжённую переменную и накопитель градиента в одну расширенную систему, интегрируемую назад. Это позволяет не хранить всю сетку прямого решателя, но заменяет память повторным вычислением траектории.

На практике существуют два рабочих режима. В первом автоматическое дифференцирование проходит через фактические шаги решателя; тогда получается точная производная дискретного алгоритма, но требуется хранить или восстанавливать его промежуточные данные. Во втором решается непрерывное сопряжённое уравнение; память может быть ниже, однако прямой и обратный проходы уже являются разными численными задачами. Выбор проверяют не по названию метода, а по времени, памяти и согласованию градиента с конечными разностями.

\begin{application}[минимальный цикл обучения Neural ODE]
Для набора пар $(x_i,y_i)$ кодировщик строит $h_i(t_0)=E_\phi(x_i)$, Neural ODE вычисляет $h_i(t_1)$, а декодер выдаёт $\widehat y_i=D_\psi(h_i(t_1))$. Потеря
\[
L(\phi,\theta,\psi)
=
\frac1N\sum_{i=1}^{N}
\ell\bigl(D_\psi(\Phi_\theta(E_\phi(x_i))),y_i\bigr)
\]
дифференцируется по параметрам кодировщика, динамики и декодера. ОДУ здесь является одним звеном обычного end-to-end конвейера, а не отдельным способом обучения.
\end{application}

\subsection{Непрерывный normalizing flow}

Пусть начальная случайная величина $z(t_0)$ имеет простую плотность $p_0$, например стандартную нормальную, и эволюционирует по уравнению
\begin{equation}
\label{eq:cnf-state-v06}
\dot z(t)=f_\theta\bigl(z(t),t\bigr).
\end{equation}
Единственность решения делает отображение $z(t_0)\mapsto z(t_1)$ обратимым. Поэтому поток переносит простое распределение в более сложное и одновременно допускает вычисление плотности.

Для дискретного обратимого преобразования $z_1=F(z_0)$ формула замены переменных содержит логарифм определителя якобиана:
\[
\log p_1(z_1)
=
\log p_0(z_0)
-
\log\left|\det D F(z_0)\right|.
\]
В непрерывном времени изменение логарифма плотности удовлетворяет скалярному ОДУ
\begin{equation}
\label{eq:instantaneous-change-of-variables-v06}
\frac{\dd}{\dd t}\log p\bigl(z(t)\bigr)
=
-\operatorname{tr}
D_zf_\theta\bigl(z(t),t\bigr).
\end{equation}
Это
\emph{мгновенная формула замены переменных}. Вместо определителя конечного преобразования используется след локального якобиана векторного поля.

\begin{definition}
\emph{Непрерывным normalizing flow}, или CNF, называется генеративная модель, в которой состояние и логарифм плотности совместно решают систему
\begin{equation}
\label{eq:cnf-augmented-v06}
\frac{\dd}{\dd t}
\begin{pmatrix}
z\\
\log p(z)
\end{pmatrix}
=
\begin{pmatrix}
f_\theta(z,t)\\
-\operatorname{tr}D_zf_\theta(z,t)
\end{pmatrix}.
\end{equation}
\end{definition}

Обучение по максимальному правдоподобию берёт объект данных $x=z(t_1)$, интегрирует поток назад к $z(t_0)$ и вычисляет
\begin{equation}
\label{eq:cnf-log-likelihood-v06}
\log p_\theta(x)
=
\log p_0\bigl(z(t_0)\bigr)
-
\int_{t_0}^{t_1}
\operatorname{tr}D_zf_\theta\bigl(z(t),t\bigr)\,\dd t.
\end{equation}
После обучения генерация идёт вперёд: сначала выбирается $z(t_0)\sim p_0$, затем решается~\eqref{eq:cnf-state-v06} до $t_1$.

След имеет простой геометрический смысл. Положительное значение $\operatorname{tr}D_zf$ означает локальное расширение объёма фазового пространства и уменьшение плотности; отрицательное --- сжатие объёма и рост плотности. CNF тем самым связывает вероятность с геометрией потока.

\begin{application}[непрерывный генератор параметров конструкции]
Пусть $z(t_0)\sim\mathcal N(0,I)$, а CNF переносит его в распределение допустимых параметров $\eta=z(t_1)$. Декодер строит по $\eta$ геометрию, а физический симулятор вычисляет характеристики конструкции. Правдоподобие обучает распределение воспроизводить примеры, а физическая цель смещает его к лёгким, прочным или устойчивым решениям. В отличие от независимой генерации и последующей фильтрации, градиент проходит через поток, декодер и симулятор как через единую модель.
\end{application}

CNF не устраняет вычислительную цену автоматически. На каждом шаге требуется значение $f_\theta$, а для плотности --- след его якобиана. В низкой размерности след можно вычислять явно. В высокой размерности необходимо использовать структуру векторного поля или приближённые оценки; иначе плотностная модель может оказаться существенно дороже обычного Neural ODE.

\subsection{Latent ODE для нерегулярных наблюдений}

Во временных рядах измерения часто приходят в разные моменты. Latent ODE задаёт непрерывную скрытую траекторию
\begin{equation}
\label{eq:latent-ode-v06}
\dot z(t)=f_\theta(z(t)),
\qquad
z(t_0)\sim q_\phi
\bigl(z(t_0)\mid x(t_1),\ldots,x(t_m)\bigr),
\end{equation}
а декодер описывает распределение наблюдения $p_\psi(x(t_i)\mid z(t_i))$. Один и тот же решатель выдаёт состояния ровно в тех временных точках, где имеются данные или нужен прогноз.

Такой подход отделяет моменты наблюдения от самой динамики. Пропуск измерения не требует искусственно заполнять ячейку регулярной сетки; состояние просто интегрируется через соответствующий интервал. Однако модель всё равно предполагает, что между наблюдениями действует одно гладкое векторное поле. Резкие переключения режимов, удары и дискретные события требуют гибридной модели, а не уменьшения допуска обычного решателя.

Для генеративного проектирования latent ODE полезен как модель скрытой эволюции испытаний. По редким измерениям деформации, температуры или вибрации кодировщик оценивает начальное скрытое состояние; ОДУ восстанавливает траекторию; декодер прогнозирует наблюдаемые характеристики в произвольные моменты. Неопределённость начального состояния отражает неоднозначность восстановления по неполным данным.

\subsection{Сквозной контур генеративного проектирования}

Рассмотрим компактную схему, объединяющую материал главы. Латентный код $\xi$ преобразуется генератором в параметры конструкции
\[
\eta=G_\psi(\xi).
\]
Они задают физическую или гибридную динамику
\begin{equation}
\label{eq:design-dynamics-v06}
\dot x
=
f_{\mathrm{phys}}(t,x,\eta)
+
r_\theta(t,x,\eta),
\qquad
x(t_0)=x_0(\eta),
\end{equation}
где $r_\theta$ корректирует неполную физическую модель по данным. Целевая функция
\begin{equation}
\label{eq:design-objective-v06}
J(\psi,\theta)
=
\Phi\bigl(x(T),\eta\bigr)
+
\int_{t_0}^{T}\ell(t,x(t),\eta)\,\dd t
+
R(\eta)
\end{equation}
может учитывать конечную деформацию, накопленную вибрацию, массу и технологические ограничения.

Прямой решатель строит траекторию. Сопряжённая система вычисляет $\nabla_\eta J$ и $\nabla_\theta J$. Затем правило цепочки передаёт сигнал генератору:
\begin{equation}
\label{eq:end-to-end-design-gradient-v06}
\nabla_\psi J
=
D_\psi G_\psi(\xi)^\top\nabla_\eta J.
\end{equation}
Если генератор является CNF, к физической цели можно добавить отрицательное лог-правдоподобие или регуляризацию распределения. Тогда одна модель одновременно описывает разнообразие проектов и их динамическое качество.

\begin{principle}[что именно должно быть дифференцируемым]
\label{pr:differentiable-design-chain-v06}
Для end-to-end обучения не требуется аналитическая формула решения ОДУ. Требуется, чтобы каждый переход в цепочке --- генератор, параметры динамики, численный решатель и целевая функция --- имел корректно вычисляемое действие производной или сопряжённого оператора.
\end{principle}

Практическая проверка такой системы проводится по уровням. Сначала на малой задаче проверяются решение и порядок сходимости решателя. Затем конечными разностями проверяется градиент физической цели по нескольким параметрам. После этого сравниваются прямой и сопряжённый градиенты, исследуется зависимость результата от допуска и только затем подключается большой генератор.

\subsection{Когда Neural ODE оправдана}

Neural ODE особенно естественна, когда время непрерывно, наблюдения нерегулярны, требуется интерполяция между ними, физическая модель уже имеет вид ОДУ или нужна обратимая деформация распределения. Её преимущества связаны со структурой задачи: адаптивным вычислением, общими параметрами динамики, непрерывной траекторией и возможностью использовать аппарат устойчивости и чувствительности.

Она менее привлекательна, когда данные статичны, фиксированная глубина работает достаточно хорошо, стоимость каждого вызова сети велика, динамика жёсткая или задача требует необратимых и разрывных преобразований. В этих случаях обычная residual-сеть, дискретная state-space модель или гибридная схема часто дают более прозрачный вычислительный бюджет.

\begin{caution}[не смешивать математическую и численную обратимость]
Точный поток при единственности решения обратим. Численная траектория при конечном допуске восстанавливается назад лишь приближённо. Поэтому обратимость модели не отменяет необходимости контролировать ошибку решателя, особенно при вычислении сопряжённого градиента и likelihood в CNF.
\end{caution}

\subsection{Итог главы}

Дифференциальное уравнение задаёт отображение не явной формулой, а траекторией. Корректность задачи Коши делает это отображение однозначным; спектр и функции Ляпунова описывают устойчивость; численный метод превращает поток в алгоритм; чувствительности и сопряжённая система дают градиент; Neural ODE помещает параметризованное векторное поле внутрь ML-модели.

Для генеративного проектирования итоговая логика имеет вид
\[
\begin{aligned}
\text{латентный код}
&\longrightarrow
\text{параметры конструкции}
\longrightarrow
\text{динамика}
\longrightarrow
\text{траектория},\\
\text{траектория}
&\longrightarrow
\text{физическая цель}
\longrightarrow
\text{сопряжённый градиент}.
\end{aligned}
\]
CNF добавляет к этой цепочке управляемое распределение проектов, а latent ODE --- работу с редкими и нерегулярными динамическими измерениями. Тем самым ОДУ становятся общим языком между физическим моделированием, обучаемой динамикой и генеративными моделями, но остаются полезными только при явном контроле корректности, устойчивости и численной ошибки.
